% ===================================================================
% PREÁMBULO — ANEXO DE CÁLCULO ESTRUCTURAL (LaTeX)
% Tesis de Grado | Edificio Mixto 18P + 3 Subsuelos | CDE, Paraguay
% Motor: pdflatex (MiKTeX) · idioma español · coma decimal
% ===================================================================
\documentclass[11pt,a4paper,oneside]{report}

% ---------- Idiomas y codificación ----------
\usepackage[utf8]{inputenc}
\usepackage[T1]{fontenc}
\usepackage[spanish,es-tabla]{babel}

% ---------- Tipografía ----------
\usepackage{lmodern}
\usepackage{microtype}

% ---------- Matemática ----------
\usepackage{amsmath}
\usepackage{amssymb}
\usepackage{mathtools}

% ---------- Unidades (SI) ----------
\usepackage{siunitx}
\sisetup{
    output-decimal-marker={,},
    exponent-product=\cdot,
    inter-unit-product=\cdot,
    detect-all,
}

% ---------- Tablas ----------
\usepackage{booktabs}
\usepackage{array}
\usepackage{multirow}
\usepackage{tabularx}
\newcolumntype{C}[1]{>{\centering\arraybackslash}p{#1}}

% ---------- Figuras ----------
\usepackage{graphicx}
\usepackage{tikz}
\usetikzlibrary{patterns,calc,positioning,babel,arrows}
\usepackage{float}
\usepackage{caption}
\captionsetup{font=small,labelfont=bf}

% ---------- Geometría de página ----------
\usepackage[a4paper,top=2.5cm,bottom=2.5cm,left=3.0cm,right=2.5cm]{geometry}

% ---------- Encabezados y pie ----------
\usepackage{fancyhdr}
\pagestyle{fancy}
\fancyhf{}
\fancyhead[L]{\small\nouppercase{\leftmark}}
\fancyhead[R]{\small\thepage}
\renewcommand{\headrulewidth}{0.4pt}

% ---------- Hipervínculos ----------
\usepackage[colorlinks=true,linkcolor=blue!60!black,
            citecolor=green!50!black,urlcolor=blue]{hyperref}

% ---------- Entornos personalizados ----------
% Recuadro "clave de diseño" (robusto con lrbox)
\newsavebox{\cajaclave}
\newenvironment{clave}[1]{%
    \par\vspace{6pt}\noindent
    \begin{lrbox}{\cajaclave}%
        \begin{minipage}{0.92\textwidth}%
            \textbf{\textcolor{blue!60!black}{#1}}\\[3pt]\small
}{%
        \end{minipage}%
    \end{lrbox}%
    \fcolorbox{blue!40}{blue!5}{\usebox{\cajaclave}}\par\vspace{6pt}%
}

% Números de ecuación con la sección
\numberwithin{equation}{section}

% Comando para nota de resultado final
\newcommand{\resultado}[1]{%
    \begin{center}
    \fbox{\parbox{0.6\textwidth}{\centering\large\textbf{#1}}}
    \end{center}}

% Comando para unidad repetida
\newcommand{\kNpm}[0]{\si{\kilo\newton\per\meter}}

\graphicspath{{figuras/}}